% =========================================
% CHAPTER 6: EXPERIMENTAL EVALUATION
% =========================================
\chapter{Experimental Results and Performance Evaluation}
\label{ch:evaluation}

This chapter presents the experimental methodology and results of the comparative evaluation. We describe the test environment and scenarios, and then analyze the results along four dimensions: throughput, batch efficiency, finality latency, and computational overhead. Each subsection includes a discussion of the observed trade-offs between the two architectures.

\section{Experimental Setup}

\subsection{Test Environment}

All experiments were conducted in a containerized environment with the following configuration:

\begin{itemize}
    \item \textbf{Hardware}: 8GB RAM, 4-core processor
    \item \textbf{OS}: Linux/Docker container environment
    \item \textbf{L1 Node}: Hardhat local node (production-like behavior)
    \item \textbf{Sequencers}: Synchronized with 5-second grace period
    \item \textbf{Duration}: Continuous 5-minute test scenarios
\end{itemize}

The Hardhat node was configured with deterministic block production (one block per transaction by default, with \texttt{evm\_mine} used to advance the chain during idle periods). Both sequencers shared this single L1 instance, ensuring that any variation in block timing or gas pricing affected them equally.

\subsection{Test Scenarios}

Three primary scenarios were defined to evaluate the framework under different load conditions:

\begin{enumerate}
    \item \textbf{Scenario~A (Normal Load)}: 150~transactions are submitted to each sequencer, representing a moderate workload intended to exercise the basic batching and submission pipeline.
    \item \textbf{Scenario~B (High Load)}: 500~transactions are submitted to stress-test the batching mechanism and measure peak throughput.
    \item \textbf{Scenario~C (Batch Analysis)}: 250~transactions are submitted to observe batching patterns and L1 submission frequency.
\end{enumerate}

In addition to these performance scenarios, two adversarial scenarios were executed:

\begin{enumerate}
    \setcounter{enumi}{3}
    \item \textbf{Scenario~D (Sequencer Crash)}: each sequencer is forcibly terminated mid-batch and restarted after a configurable delay, measuring recovery time and transaction loss.
    \item \textbf{Scenario~E (Invalid State Proposal)}: the sequencer submits a deliberately malformed state root, testing the detection mechanism of each verifier contract.
\end{enumerate}

The reproducible measurements reported in this chapter are limited to the raw artefact \texttt{measurement/raw/baseline-path-a.json}: batch-commitment and mock-proof timings, per-submission gas, and the deterministic crash-loss invariant. Screenshots of the live dashboard and Grafana are included as illustrations of the running system; they are not measurement evidence. TPS values displayed by the demo are dominated by the rate at which the harness injects HTTP requests and are not treated as an architectural throughput result.

\subsection{Data Validation Before Analysis}

Before interpreting screenshots, two checks are applied. First, the transaction counters in the Flask dashboard must match the number of transactions submitted, so that a dropped request or a metrics lag is visible. Second, Prometheus must show both sequencer scrape targets as up; zero values after a completed scenario indicate a scrape gap and require a longer post-completion wait. The \texttt{*\_finality\_time\_seconds} histograms record L1 submission-to-receipt latency (sub-second in Hardhat), not protocol finality; no claim about finality is derived from them.

These checks apply to the live-system screenshots in Figures~\ref{fig:flask-dashboard}, \ref{fig:grafana-comparison}, \ref{fig:throughput-comparison}, and~\ref{fig:finality-comparison}. The latter two are Grafana panel screenshots illustrating the dashboard, not measured throughput or finality results; no numeric claim is based on them.

Figure~\ref{fig:flask-dashboard} shows the Flask real-time dashboard with both sequencers running. Figure~\ref{fig:grafana-comparison} shows the Grafana metrics dashboard with all comparison panels visible.

\begin{figure}[htbp]
    \centering
    \includegraphics[width=0.95\textwidth]{dashboards/flask-dashboard-overview.png}
    \caption{Flask real-time dashboard during a live run. Both panels poll their sequencer's \texttt{/metrics} endpoint; the counters and batch totals depend on which scenarios were run in that container and are illustrative, not a reproducible measurement. The Architectural Comparison table records design-level properties only and is marked ``not evidence'' in the UI.}
    \label{fig:flask-dashboard}
\end{figure}

\begin{figure}[htbp]
    \centering
    \includegraphics[width=0.95\textwidth]{dashboards/grafana-top-panels.png}
    \caption{Grafana dashboard top row: transaction counters, batch counts, current TPS and L1 submission-to-receipt latency (P95) for both architectures. The latency panel measures submission latency only; protocol finality (5-block and 10-block rules in the demo) is enforced by the contracts and is not visible here. No numeric latency value is quoted from this screenshot.}
    \label{fig:grafana-comparison}
\end{figure}

\section{Throughput Analysis}

\subsection{Transactions Per Second}

Table~\ref{tab:tps-results} records the TPS gauges displayed by the two sequencers for the three demo scenarios. These readings are shown only to illustrate the live dashboard: they are dominated by the rate at which the harness injects HTTP requests and are not an architectural throughput measurement. No percentage advantage is claimed from them.

\begin{table}[h]
    \centering
    \begin{tabular}{l|c|c}
        \hline
        \textbf{Scenario} & \textbf{ZK gauge} & \textbf{Optimistic gauge} \\
        \hline
        Normal (150~tx) & illustrative & illustrative \\
        \hline
        High Load (500~tx) & illustrative & illustrative \\
        \hline
        Batch (250~tx) & illustrative & illustrative \\
        \hline
    \end{tabular}
    \caption{Illustrative TPS gauge readings from the live demo; not a throughput benchmark and not measurement evidence.}
    \label{tab:tps-results}
\end{table}

The two sequencers differ in configuration --- a 100-transaction batch every 5 seconds for the ZK path and a 200-transaction batch every 8 seconds for the Optimistic path --- which changes how frequently batches are submitted. Any rate difference observed in the demo follows from these settings and the injection harness, not from the proof mechanism.

\begin{figure}[htbp]
    \centering
    \includegraphics[width=0.95\textwidth]{dashboards/grafana-bottom-panels.png}
    \caption{Grafana dashboard lower panels showing the transaction rate comparison, ZK mock proof generation time, Optimistic batch commitment computation time, pending transaction queues, and gas per batch. These panels illustrate live metrics; the reproducible gas figures are reported in Section~\ref{sec:gas}.}
    \label{fig:throughput-comparison}
\end{figure}

\subsection{Batch Efficiency}

The Optimistic path makes more efficient use of each L1 submission because its configured batch size is twice as large. Table~\ref{tab:batch-efficiency} records this configuration-derived relationship, not a measured run.

\begin{table}[h]
    \centering
    \begin{tabular}{l|c|c|c}
        \hline
        \textbf{Metric} & \textbf{ZK Rollup} & \textbf{Optimistic} & \textbf{Trade-off} \\
        \hline
        Batch threshold & 100~txs & 200~txs & 2$\times$ difference \\
        \hline
        Batches per 200~tx & 2 & 1 & 50\% fewer for Optimistic \\
        \hline
        L1 submission frequency & higher & lower & L1 cost versus latency \\
        \hline
    \end{tabular}
    \caption{Batch efficiency derived from the configured thresholds (not a measured workload).}
    \label{tab:batch-efficiency}
\end{table}

We define batch efficiency as the ratio of transactions processed to L1 submissions:
\begin{equation}
\text{Batch Efficiency} = \frac{\text{Transactions Batched}}{\text{L1 Submissions}}
\end{equation}

For $n$ transactions the ZK path submits $\lceil n/100 \rceil$ batches and the Optimistic path $\lceil n/200 \rceil$, so the latter makes approximately half as many L1 submissions for the same workload. This follows directly from the configuration, not from a measured performance difference. It illustrates the trade-off: the ZK path favours earlier submission (smaller batches), while the Optimistic path favours fewer on-chain transactions (larger batches).

\section{Finality Analysis}

\subsection{L1 Submission Latency}

The sequencers export a histogram of the time from \texttt{send\_raw\_transaction} to receipt. Hardhat's automining mode mines each transaction immediately, so these observations reflect Python RPC overhead rather than block production wait time. No raw artefact for this histogram is committed; the table therefore records its scope rather than a P95 value.

\begin{table}[h]
    \centering
    \begin{tabular}{l|c|c}
        \hline
        \textbf{Metric} & \textbf{ZK Rollup} & \textbf{Optimistic} \\
        \hline
        What is measured & submission-to-receipt & submission-to-receipt \\
        \hline
        Typical magnitude in Hardhat & sub-second & sub-second \\
        \hline
        Protocol finality included? & No & No \\
        \hline
        Committed raw artefact? & No & No \\
        \hline
    \end{tabular}
    \caption{Scope of the \texttt{*\_finality\_time\_seconds} histograms; no latency figure is quoted from them.}
    \label{tab:l1-latency}
\end{table}

Both sequencers complete L1 submission in well under a second in Hardhat's automining mode. The Grafana panels formerly labelled ``Finality Time'' now use the name \emph{L1 Submission Latency} to make this distinction explicit. This chapter does not quote a P95 value from a screenshot. The finality difference is enforced by contract rules, not by the submission step.

\subsection{Finality Rules (Not Measured Latencies)}

L1 inclusion is a necessary but not sufficient condition for finality. Each architecture imposes an additional rule before a transaction is considered final. Table~\ref{tab:true-finality} describes those rules; the wall-clock figures sometimes quoted in prose (about 15 seconds for ZK and seven days for Optimistic) are protocol assumptions, not measurements of this stack.

\begin{table}[h]
    \centering
    \begin{tabular}{l|c|c}
        \hline
        \textbf{Finality} & \textbf{ZK Rollup} & \textbf{Optimistic} \\
        \hline
        Rule in this testbed & proof accepted + 5 L1 blocks & 10 L1 blocks (challenge window) \\
        \hline
        Production assumption & $\sim$15~s & $\sim$7~days \\
        \hline
        Hardhat wall-clock duration & none (mines on transactions only) & none (mines on transactions only) \\
        \hline
    \end{tabular}
    \caption{Finality rules and assumptions; these are not measured latencies.}
    \label{tab:true-finality}
\end{table}

The ZK rule is architecturally faster than the Optimistic rule because it does not wait for a challenge window. In production this is commonly summarised as roughly 15 seconds versus a seven-day window, a ratio of about 40,320$\times$ (604,800\,s / 15\,s). That figure is an assumption about protocol parameters, not an experimental result; the earlier 672$\times$ figure was an arithmetic error. Hardhat mines only on transactions, so the demo has no meaningful wall-clock finality, and the previously quoted 8$\times$ demo ratio was a harness artefact. The practical implication for latency-sensitive applications follows from the rule, not from a measured timing.

\begin{figure}[htbp]
    \centering
    \includegraphics[width=0.95\textwidth]{dashboards/grafana-top-panels.png}
    \caption{The same Grafana top-row screenshot as Figure~\ref{fig:grafana-comparison}, repeated here next to the finality discussion. Its panels show L1 submission-to-receipt latency, not protocol finality; the 5-block and 10-block rules are enforced by the contracts and are not visible in this metric. No latency value should be cited from the image.}
    \label{fig:finality-comparison}
\end{figure}

\section{Computational Overhead: Proof Generation}

Both sequencers first compute the same canonical keccak256 binary Merkle commitment over ABI-encoded \texttt{(from,to,value)} tuples via \texttt{l2\_common.commitment.batch\_merkle\_root}. The ZK path additionally computes a mock proof. Table~\ref{tab:proof-latency} reports the reproducible Path A measurements from \texttt{measurement/raw/baseline-path-a.json}.

\begin{table}[h]
    \centering
    \begin{tabular}{l|c|c}
        \hline
        \textbf{Metric} & \textbf{ZK path} & \textbf{Optimistic path} \\
        \hline
        Batch commitment P95 (100~tx) & 3.64\,ms & 3.64\,ms \\
        \hline
        Batch commitment P95 (200~tx) & 6.80\,ms & 6.80\,ms \\
        \hline
        Mock proof P95 (any batch size) & 0.008\,ms & --- (no proof step) \\
        \hline
    \end{tabular}
    \caption{Reproducible Path A timings (P95); source: \texttt{measurement/raw/baseline-path-a.json}.}
    \label{tab:proof-latency}
\end{table}

The mock proof is a single domain-separated keccak hash of the commitment --- not a SNARK --- taking roughly 0.008\,ms at P95. Both paths therefore pay the same commitment cost, and the ZK-only proof step is about three orders of magnitude smaller. This is a lower bound specific to the mock: a real Groth16 prover would take seconds, so the figure cannot be extrapolated to real prover hardware. The plan for such a prover is documented in \texttt{docs/PATH\_B\_GROTH16.md} and is future work, not part of this evaluation.

\section{Security Analysis}
\label{sec:security}

The performance results above have direct security implications. This section examines the threat model, compares the two architectures' security properties, and presents the empirical results of the two adversarial scenarios that could be tested directly.

\subsection{Threat Model}

Five adversarial scenarios are relevant to any Layer~2 deployment. We assume the adversary controls the sequencer process but not the Layer~1 consensus mechanism, and that at least two-thirds of L1 validators are honest.

\begin{enumerate}
    \item \textbf{Malicious sequencer}: submits an invalid state root to L1.
    \item \textbf{L1 reorganization}: confirmed L1 blocks are replaced, invalidating anchored commitments.
    \item \textbf{Sequencer crash}: the sequencer fails mid-batch, losing in-flight transactions.
    \item \textbf{Transaction censorship}: the sequencer selectively excludes certain transactions.
    \item \textbf{Data unavailability}: the sequencer withholds L2 batch data, preventing independent state reconstruction.
\end{enumerate}

Scenarios~1 and 3 were directly tested (Sections~\ref{sec:invalid-state} and~\ref{sec:crash} respectively). Scenarios~2, 4, and 5 are analyzed architecturally; Hardhat's deterministic block production does not simulate chain reorganizations, network-level censorship, or data withholding.

\subsection{Security Properties Compared}

Table~\ref{tab:security-models} contrasts the key security properties of each architecture.

\begin{table}[h]
    \centering
    \begin{tabular}{l|>{\raggedright\arraybackslash}p{4cm}|>{\raggedright\arraybackslash}p{4cm}}
        \hline
        \textbf{Property} & \textbf{ZK Rollup} & \textbf{Optimistic Rollup} \\
        \hline
        Security assumption & 1-of-N honest (one prover suffices) & At least one honest challenger required \\
        \hline
        Proof type & Cryptographic validity proof & Economic fraud proof \\
        \hline
        Invalid state halted & Immediately upon proof rejection & After fraud proof within challenge window \\
        \hline
        Cost model & Computational (proof generation) & Economic (challenge period and bonds) \\
        \hline
        Censorship risk & Low (anyone can generate a proof) & Medium (requires active challengers) \\
        \hline
    \end{tabular}
    \caption{Security model comparison. These are properties of the architectures in principle; in this testbed the ZK proof is a mock and the on-chain verifier performs only a non-zero check, so none of the guarantees listed here is implemented.}
    \label{tab:security-models}
\end{table}

Architecturally, the ZK model rests on the soundness of a proof system: a real validity proof would guarantee the state transition, while the Optimistic model rests on economic incentives and at least one honest challenger. In this testbed the ZK proof is a domain-separated keccak mock and \texttt{ZKVerifier.verifyProof} checks only non-zero words; the \emph{interface} is preserved but the cryptographic guarantee is not implemented. The distinction is documented in \texttt{docs/IMPLEMENTATION\_NOTES.md}.

\subsection{Challenger Absence: Quantifying the Risk}
\label{sec:challenger-absence}

The most significant vulnerability in an Optimistic~Rollup is that no challenger submits a fraud proof before the window closes. If each of $n$ independent observers detects fraud with probability $p$, the probability that \emph{none} of them challenges is:

\begin{equation}
P(\text{no challenge}) = (1 - p)^n
\end{equation}

With $p = 0.1$ and $n = 10$ observers, $P(\text{no challenge}) = (0.9)^{10} \approx 0.35$---a 35\% failure rate even with ten watchers. To bring this below a target threshold $\varepsilon$, the minimum observer count required is:

\begin{equation}
n \geq \frac{\ln \varepsilon}{\ln (1 - p)}
\end{equation}

For $p = 0.1$ and $\varepsilon = 0.01$, this gives $n \geq 44$~independent watchtower nodes. This is an illustration of the monitoring burden implied by the optimistic assumption, not a measurement of this testbed. A production ZK rollup has no corresponding watchtower requirement because a real validity proof rejects invalid state; the mock verifier here does not implement that guarantee.

\subsection{Behavior Under Invalid State Proposals}
\label{sec:invalid-state}

Scenario~E exercised malicious-root handling in the mocked adversarial harness. The Python tests configure a verifier \texttt{MagicMock} to return \texttt{False} and assert that the sequencer treats the commitment as not finalised. This tests the sequencer's control flow only: no on-chain rejection is exercised. The deployed \texttt{ZKVerifier} accepts any non-zero 32-byte proof word and has no binding between proof and root, so an on-chain mock verifier would \emph{not} reject a mismatched proof.

Architecturally, a real ZK rollup would reject an invalid state at proof verification (\emph{prevent-by-default}), while an optimistic rollup accepts it provisionally and relies on a challenge within the window (\emph{detect-and-revert}). In this testbed only the optimistic status transition (marking a commitment \texttt{CHALLENGED}) is modelled; no real fraud-proof submission is implemented, and no wall-clock timing for a challenge is measured.

\section{Gas Cost Analysis}
\label{sec:gas}

\emph{Important scope note:} The gas figures in this section come from a local Hardhat node with a fixed gas price and deterministic block production; they are useful only for a \emph{relative} comparison, not as production cost estimates. In this testbed the Optimistic submission costs slightly \emph{more} gas per transaction (170,937 versus 165,293) because its commitment struct writes more storage fields; the mock proof adds no meaningful calldata cost. A real Groth16 verifier would change this picture substantially, but that is future work (see \texttt{docs/PATH\_B\_GROTH16.md}), not a measured result.

\subsection{Per-Submission Gas (Hardhat measurements)}

Table~\ref{tab:gas-costs} reports the gas consumed per L1 submission measured directly from transaction receipts via \texttt{tx.gasUsed} in the Hardhat environment (see Appendix~\ref{app:metrics} for raw values).

\begin{table}[h]
    \centering
    \begin{tabular}{l|c|c}
        \hline
        \textbf{Metric} & \textbf{ZK Rollup} & \textbf{Optimistic} \\
        \hline
        Gas per submission (measured) & 165,293 & 170,937 \\
        \hline
        Batch threshold & 100 tx & 200 tx \\
        \hline
        Source & Hardhat receipt & Hardhat receipt \\
        \hline
    \end{tabular}
    \caption{L1 gas per submission measured from Hardhat receipts; no multi-batch total is reported because no raw artefact supports one.}
    \label{tab:gas-costs}
\end{table}

The per-submission gas costs are close (165,293 for ZK versus 170,937 for Optimistic). The Optimistic contract is slightly more expensive because its \texttt{Commitment} struct writes more storage fields. These are single-submission local measurements supporting only a relative comparison; the earlier 1.52$\times$ total-cost and 11-versus-7-batch figures had no reproducible artefact and were removed. Any estimate of real-Groth16 verifier gas is future-work speculation, not a measurement (see \texttt{docs/PATH\_B\_GROTH16.md}).

\subsection{Per-Transaction Cost}

Amortising the per-submission cost over a full batch gives an indicative figure:

\begin{equation}
\text{Cost}_{\text{ZK}} \approx \frac{165{,}293}{100} \approx 1{,}653 \text{ gas/tx}, \quad
\text{Cost}_{\text{Opt}} \approx \frac{170{,}937}{200} \approx 855 \text{ gas/tx}
\end{equation}

These figures assume full batches and follow from the configured thresholds, not from a measured workload. The Optimistic path is cheaper per transaction because it submits half as often; translating either figure into a dollar cost requires a real mainnet gas price, which Hardhat does not provide.

\section{Adversarial Scenario Results}
\label{sec:scenarios}

\subsection{Sequencer Crash and Recovery (Scenario~D)}
\label{sec:crash}

To evaluate resilience, each sequencer was forcibly terminated (\texttt{SIGKILL}) while it had pending transactions in its pool, then restarted after a 10-second delay. Table~\ref{tab:crash-recovery} reports the recovery behavior.

\begin{table}[h]
    \centering
    \begin{tabular}{l|c|c}
        \hline
        \textbf{Metric} & \textbf{ZK Rollup} & \textbf{Optimistic} \\
        \hline
        Accepted but unconfirmed at crash & 100 & 200 \\
        \hline
        Deterministic loss (batch threshold) & 100 & 200 \\
        \hline
        Pool persistence across restart & none & none \\
        \hline
        First batch after restart & after next threshold & after next threshold \\
        \hline
        L1 state consistency preserved & Yes & Yes \\
        \hline
    \end{tabular}
    \caption{Crash-loss invariant measured with the real sequencer classes under mocked RPC: the accepted-but-unconfirmed set equals the batch threshold.}
    \label{tab:crash-recovery}
\end{table}

Both sequencers hold their transaction pool in memory only, so a crash loses the accepted-but-unconfirmed set. The reproducible invariant, recorded in \texttt{measurement/raw/baseline-path-a.json}, is deterministic: 100 transactions for ZK and 200 for Optimistic, matching the respective batch thresholds. The earlier 73/142 figures had no raw artefact and were removed. L1 state integrity was preserved in both cases: no partial or inconsistent batches were committed.

The larger loss window on the Optimistic path follows directly from its larger batch threshold. This is a design trade-off: larger batches reduce on-chain submission frequency but increase the number of transactions exposed to a sequencer failure.

\section{Discussion}

\subsection{Answering the Research Questions}

We now revisit the research questions posed in Chapter~\ref{ch:intro} and evaluate the extent to which our experimental results provide answers.

\textbf{RQ1 (Finality).} Finality is a protocol rule, not a measured wall-clock value. The ZK rule accepts a proof after 5 L1 blocks; the Optimistic rule waits 10 L1 blocks in the demo (about seven days in production). Assuming roughly 15 seconds for ZK versus seven days for Optimistic gives an assumed gap of about 40,320$\times$. Both architectures share the same L1 submission step. This is an architectural comparison, not a measured latency result.

\textbf{RQ2 (L1 Congestion and Reorgs).} Both architectures are equally affected by L1 block production delays at the submission step. A reorganisation analysis can only be architectural here: Hardhat's deterministic block production does not simulate reorgs. For a real ZK rollup the recovery path would be detect, re-prove and re-include; in this testbed the mock proof takes about 0.008\,ms (P95), so the re-proving term is negligible by construction. An Optimistic rollup must re-run the dispute window on the new canonical chain --- potentially seven days in production. No wall-clock reorg recovery time was measured.

\textbf{RQ3 (Censorship Resistance).} The crash-loss invariant (Table~\ref{tab:crash-recovery}) shows that both sequencers preserve L1 state integrity but lose their in-memory pool: 100 transactions for ZK and 200 for Optimistic, matching their batch thresholds. The larger exposure on the Optimistic path follows from its larger batch size; neither path implements a persistent mempool.

\textbf{RQ4 (Fault-Escape Mechanisms).} Architecturally the ZK model rejects invalid state via proof verification while the Optimistic model relies on an active challenger. In this testbed the ZK verifier is a mock that accepts any non-zero proof word, and the adversarial test exercises a mocked verifier returning \texttt{False}; no on-chain rejection is demonstrated. The economic asymmetry is a property of the architecture, not a measured result here.

\textbf{RQ5 (Data Availability).} Both sequencers submit the batch commitment to L1; the full transaction data backing that commitment is not posted in this prototype, so the framework does not demonstrate data availability. Architecturally, posting batch data on-chain would let observers reconstruct state; the gas implications are discussed in Section~\ref{sec:gas}.

\subsection{Summary of Trade-offs}

Table~\ref{tab:tradeoff-summary} consolidates the findings across all evaluation dimensions. Measured entries come from \texttt{measurement/raw/baseline-path-a.json}; architectural entries are consequences of the design rather than measured results.

\begin{table}[h]
    \centering
    \begin{tabular}{l|c|c|l}
        \hline
        \textbf{Dimension} & \textbf{ZK Rollup} & \textbf{Optimistic} & \textbf{Basis} \\
        \hline
        Batch commitment P95 & 3.64\,ms (100~tx) & 3.64\,ms (100~tx) & measured \\
        \hline
        Mock proof P95 & 0.008\,ms & none & measured \\
        \hline
        Batch threshold & 100~tx & 200~tx & configuration \\
        \hline
        Gas per submission & 165,293 & 170,937 & measured \\
        \hline
        Finality rule & proof + 5 L1 blocks & 10 L1 blocks & protocol rule \\
        \hline
        Crash loss (accepted, unconfirmed) & 100 & 200 & measured \\
        \hline
        Validity guarantee & mock verifier & challenge model & not implemented \\
        \hline
    \end{tabular}
    \caption{Consolidated trade-off summary. Measured values cite the Path A raw artefact; no unsupported performance claims are included.}
    \label{tab:tradeoff-summary}
\end{table}

Neither architecture dominates across all dimensions. The ZK path finalises after a shorter rule and has a smaller crash-loss window; the Optimistic path submits less often and is cheaper per transaction in this local testbed. The proof guarantee itself is a mock, so the security comparison is architectural. The choice between them remains application-dependent.

\section{Threats to Validity}

Several factors may limit the generalizability of our experimental results.

\textbf{Internal validity.} The Hardhat local node does not replicate the full complexity of Ethereum mainnet. Block production is deterministic rather than stochastic, gas pricing is fixed rather than market-driven, and there is no network latency between the sequencer and the L1 node. These simplifications mean that absolute latency and gas values should be interpreted as relative comparisons between the two architectures rather than production-grade predictions.

\textbf{Construct validity.} The ZK proof is a single domain-separated keccak hash (about 0.008\,ms at P95), not a SNARK. It preserves the structural pattern of the pipeline (commit, prove, verify) but vastly underestimates real proving cost; no overhead percentage is claimed from it. A real Groth16 prover is documented as future work in \texttt{docs/PATH\_B\_GROTH16.md}.

\textbf{External validity.} The framework evaluates only two rollup architectures (ZK and Optimistic) with a single configuration each. Varying the batch size, challenge period, or proof system could yield different trade-off profiles. Our results are specific to the chosen parameter settings and should not be extrapolated without re-evaluation under alternative configurations.

\textbf{Statistical validity.} The measurements reported here are P95 values from the reproducible Path A harness recorded in a single raw artefact, not means over repeated experimental runs. No confidence interval or hypothesis test is claimed; a larger-scale study would be needed to support one.

\section{Summary}

Three findings stand out from this chapter. First, the finality gap is architectural: the ZK rule accepts a proof plus five L1 blocks, while the Optimistic rule waits a challenge window of ten L1 blocks in the demo (about seven days in production). Assuming 15 seconds for ZK gives an assumed ratio of roughly 40,320$\times$, not a measurement. Second, per-submission gas is similar (165,293 versus 170,937 on Hardhat), with the Optimistic path slightly more expensive because of its larger commitment struct; no multi-batch total is reported. Third, crash loss is deterministic at 100 ZK versus 200 Optimistic transactions, matching the batch thresholds.

The security analysis (Section~\ref{sec:security}) describes architectural posture: a real ZK rollup prevents invalid state by construction, while an optimistic rollup needs active challengers and the illustrative 44-node monitoring calculation. In this prototype the ZK verifier is a mock, so the guarantee is an interface, not an implemented cryptographic property.
